\documentclass[11pt]{article}
\usepackage[a4paper,margin=2.8cm]{geometry}
\usepackage{amsmath,amssymb,amsthm}
\usepackage{hyperref}

\newtheorem{definition}{Definition}
\newtheorem{lemma}{Lemma}
\newtheorem{theorem}{Theorem}
\theoremstyle{remark}
\newtheorem{remark}{Remark}

\newcommand{\Zar}{\operatorname{Zar}}
\newcommand{\Spec}{\operatorname{Spec}}
\newcommand{\cO}{\mathcal{O}}
\newcommand{\kw}[1]{\texttt{#1}}

% stamp on every page: file name and date/time of compilation
\newcount\stamphour \newcount\stampmin
\stamphour=\time \divide\stamphour by 60
\stampmin=\time \multiply\stamphour by 60 \advance\stampmin by -\stamphour \divide\stamphour by 60
\edef\stamptime{\the\year-\ifnum\month<10 0\fi\the\month-\ifnum\day<10 0\fi\the\day\space
 \ifnum\stamphour<10 0\fi\the\stamphour:\ifnum\stampmin<10 0\fi\the\stampmin}
\AddToHook{shipout/foreground}{\put(0.5\paperwidth,-\dimexpr\paperheight-0.45in\relax)%
 {\makebox[0pt]{\scriptsize\ttfamily\jobname.pdf\quad\stamptime}}}

\title{The line with doubled origin is not separated}
\author{}
\date{}

\begin{document}
\maketitle

Rings are commutative. Schemes are represented, as in~\cite{CLS,Zeu}, by
locally ringed lattices: a distributive lattice $L$ of (compact) opens with
a sheaf of rings $\cO$ on it. For a ring $A$, $\Spec A$ is the Zariski
lattice $\Zar(A)$, generated by the basic opens $D(f)$, with the structure
sheaf $\cO_A$; we use only that $\cO_A(1)=A$, that
$\cO_A(D(f))=A[1/f]$, and that the restriction $\cO_A(1)\to\cO_A(D(f))$
is the canonical map $A\to A[1/f]$. A locally ringed lattice $X$ is
\emph{affine} when the canonical map $X\to\Spec\Gamma(X)$, with
$\Gamma(X)=\cO(1)$, is an isomorphism, and an open $u$ is \emph{affine}
when the restriction $X|_u$ is affine.

\section{Separated schemes}

For opens $u,v$ of $X$ write $r\colon\cO(u)\to\cO(u\wedge v)$ and
$r'\colon\cO(v)\to\cO(u\wedge v)$ for the restriction maps.

\begin{definition}
The overlap $u\wedge v$ is \emph{generated} by $u$ and $v$ when every
$z\in\cO(u\wedge v)$ is a finite sum
\[
  z=r(a_1)\,r'(b_1)+\dots+r(a_n)\,r'(b_n),\qquad a_i\in\cO(u),\ b_i\in\cO(v),
\]
that is, when $\cO(u)\otimes\cO(v)\to\cO(u\wedge v)$ is surjective. $X$
is \emph{separated} when for any two affine opens $u$ and $v$ the overlap
$u\wedge v$ is affine and generated by $u$ and $v$.
\end{definition}

\begin{remark}
Classically, a scheme is separated, i.e.\ its diagonal is a closed
immersion, exactly when this holds for all affine opens $u,v$: the
diagonal restricted to $\Spec\cO(u)\times\Spec\cO(v)$ is the closed
immersion $\Spec\cO(u\wedge v)\to\Spec(\cO(u)\otimes\cO(v))$. The
definition above avoids covers, products of schemes and tensor products.
\end{remark}

\section{The line with doubled origin}

Let $k$ be a ring, $A=k[X]$ and $d=D(X)\in\Zar(A)$, the punctured line.
The line with doubled origin glues two copies of $\Spec A$ along $d$ by
the identity. Its lattice is
\[
  L=\{(u_0,u_1)\in\Zar(A)\times\Zar(A)\mid u_0\wedge d=u_1\wedge d\},
\]
with componentwise order and operations, and its sheaf is
\[
  \cO(u_0,u_1)=\cO_A(u_0)\times_{\cO_A(m)}\cO_A(u_1),\qquad
  m=u_0\wedge d=u_1\wedge d,
\]
the pairs $(s,t)$ with $s|_m=t|_m$, restricted componentwise. The two
charts are
\[
  e_0=(1,d),\qquad e_1=(d,1),
\]
the two copies of the line; each is isomorphic to $\Spec A$, hence
affine. Their overlap is $e_0\wedge e_1=(d,d)$, and
$\cO(e_0\wedge e_1)$ consists of the pairs $(s,t)$ of elements of
$\cO_A(d)=A[1/X]$ with $s=t$.

\begin{lemma}\label{lem:lift}
Let $\rho\colon A\to C$ and $\pi\colon B\to C$ be ring maps. The elements
$z\in B$ such that $\pi(z)$ is in the image of $\rho$ form a subring of
$B$. In particular they are closed under finite sums of products.
\end{lemma}

\begin{proof}
The image of a ring map is a subring, and $\pi^{-1}$ of a subring is a
subring.
\end{proof}

\begin{theorem}
If the line with doubled origin over $k$ is separated, then $1=0$ in $k$.
\end{theorem}

\begin{proof}
Assume it is separated. The charts $e_0$ and $e_1$ are affine, so their
overlap $(d,d)$ is generated by them. Let $\rho\colon A\to A[1/X]$ be the
canonical map and $\pi\colon\cO(d,d)\to A[1/X]$ the first projection.

\emph{Sections of $e_0$.} An element of $\cO(e_0)$ is a pair $(a_0,a_1)$
with $a_0\in\cO_A(1)=A$. The first component of its restriction to
$(d,d)$ is $a_0|_d=\rho(a_0)$.

\emph{Sections of $e_1$.} An element of $\cO(e_1)$ is a pair $(b_0,b_1)$
with $b_0\in\cO_A(d)$, $b_1\in\cO_A(1)=A$, and $b_0|_m=b_1|_m$ for
$m=d\wedge d=d$. So $b_0=b_1|_d$, and the first component of the
restriction of $(b_0,b_1)$ to $(d,d)$ is $\rho(b_1)$.

Hence $\pi$ maps the restrictions of both $\cO(e_0)$ and $\cO(e_1)$ into
the image of $\rho$, and by Lemma~\ref{lem:lift} so does it map every
finite sum of products $r(a_1)r'(b_1)+\dots+r(a_n)r'(b_n)$.

Now take $z_0=(1/X,1/X)\in\cO(d,d)$. By generation, $1/X=\rho(p)$ for
some polynomial $p\in k[X]$. Then $\rho(Xp)=X\cdot\frac1X=1=\rho(1)$. The
map $\rho$ is injective, since $X$ is a regular element of $k[X]$: if
$q/1=0$ in $A[1/X]$ then $X^nq=0$ for some $n$, so $q=0$. Hence $Xp=1$ in
$k[X]$, and comparing constant coefficients gives $0=1$ in $k$.
\end{proof}

Only the generation condition is used, and only at one pair of affine
opens. The overlap $(d,d)$ is in fact affine, isomorphic to
$\Spec k[X,X^{-1}]$, so it is the second condition that fails: the two
copies of $k[X]$ both map onto the same subring $k[X]$ of $k[X,X^{-1}]$,
which does not contain $1/X$. Generation is a proposition (the sums are
merely given), and the conclusion $1=0$ is a proposition since $k$ is a
set, so the argument is constructive as written. In particular, over a
ring with $1\neq0$ the doubled line is not separated.

\section{What has been formalised}

The definition and the theorem are checked in the proof checker
\textsc{Unit}~\cite{Unit}, in the development \texttt{sag/qcqs} of qcqs
schemes as locally ringed lattices.
\begin{itemize}
\item \texttt{SchemeSeparated.seg} defines \kw{Overlap[u,v].Generated}
  (every section of the overlap is merely a finite sum of products of
  restrictions) and \kw{IsSeparated}, with ``affine open'' the existing
  notion \kw{Qcqs.AffineOpen}.
\item \texttt{DoubledLineNotSeparated.seg} proves
  \[
    \kw{separated-one-zero} : \kw{IsSeparated}\to 1=0
  \]
  for the doubled line \kw{Dbl.doubled} of \texttt{DoubledLine.seg}, over
  any ring $k$. It follows the proof above: the lifting
  lemma, the two chart computations (the second uses the agreement of the
  components on the middle open), the element $z_0$ built from the
  inverse of $X$ in the localisation, and injectivity of
  $k[X]\to k[X,1/X]$ with evaluation at $0$
  (\texttt{VariableKernel.seg}). The affineness of the charts is
  \texttt{DoubledLineQcqs.seg}.
\item Not formalised: that the overlap $(d,d)$ is affine (not needed),
  and the classical equivalence of the definition with the diagonal
  being a closed immersion (the definition is taken as primitive).
\end{itemize}
The development \texttt{sag/qcqs} checks with 16\,141 declarations and no
failure, in about 19\,s. No new axiom is used; the parameters of the
development are univalence, images and surjections, as for the rest of
\texttt{sag/qcqs}. The doubled line is also checked to be a qcqs scheme
(\texttt{DoubledLineQcqs.seg}) and not affine over a nonzero ring
(\texttt{DoubledLineNotAffine.seg}).

\begin{thebibliography}{9}
\bibitem{CLS} T.~Coquand, H.~Lombardi and P.~Schuster. Spectral schemes
  as ringed lattices. \emph{Annals of Mathematics and Artificial
  Intelligence} 56 (2009), 339--360.
\bibitem{Unit} T.~Coquand. \textsc{Unit}: a type checker for dependent
  type theory. \url{https://github.com/coquand/unit}.
\bibitem{Zeu} M.~Zeuner. Univalent foundations of constructive
  algebraic geometry. Stockholm University, 2024.
\end{thebibliography}

\end{document}
